% ============================================================================
% ARBEITSBLATT: Bohr-Vertiefung (LP-07) - Klasse 12 GK
% Stand: 06.02.2026
% ============================================================================

\documentclass[11pt,a4paper]{article}

\usepackage[utf8]{inputenc}
\usepackage[T1]{fontenc}
\usepackage[ngerman]{babel}
\usepackage{geometry}
\usepackage{helvet}
\usepackage{fancyhdr}
\usepackage{xcolor}
\usepackage{amsmath,amssymb}
\usepackage{siunitx}
\sisetup{locale = DE}
\usepackage{tcolorbox}
\usepackage{array}
\usepackage{tabularx}
\usepackage{booktabs}
\usepackage{tikz}
\usepackage{qrcode}
\usepackage[hidelinks]{hyperref}

\renewcommand{\familydefault}{\sfdefault}
\geometry{left=1.5cm, right=1.5cm, top=1.2cm, bottom=1.0cm}
\setlength{\parindent}{0pt}
\setlength{\parskip}{0pt}

% Farben
\definecolor{physik}{HTML}{2c5aa0}
\definecolor{hinweis}{HTML}{fff3e0}
\definecolor{beispiel}{HTML}{e3f2fd}

% Kopfzeile
\pagestyle{fancy}
\fancyhf{}
\fancyhead[L]{\small\textbf{Physik 12 GK} | Bohr-Vertiefung}
\fancyhead[R]{\small Arbeitsblatt}
\renewcommand{\headrulewidth}{0.4pt}
\setlength{\headheight}{14pt}

% Boxen
\tcbuselibrary{skins}
\newtcolorbox{merkbox}[1][]{colback=white, colframe=physik, boxrule=1.5pt, arc=0pt, left=6pt, right=6pt, top=6pt, bottom=6pt, #1}
\newtcolorbox{formelbox}[1][]{colback=beispiel, colframe=physik, boxrule=1pt, arc=0pt, left=6pt, right=6pt, top=4pt, bottom=4pt, #1}
\newtcolorbox{hinweisbox}[1][]{colback=hinweis, colframe=orange!70!black, boxrule=1pt, arc=0pt, left=6pt, right=6pt, top=4pt, bottom=4pt, #1}

% Befehle
\newcommand{\aufgabe}[1]{\textbf{\textcolor{physik}{#1}}}
\newcommand{\luecke}[1]{\underline{\hspace{#1}}}
\newcommand{\niveau}[1]{\textcolor{physik}{\small (#1)}}
\newcommand{\sterne}[1]{\niveau{\ifcase#1\or$\star$\or$\star\star$\or$\star\star\star$\fi}}

\begin{document}

% ============================================================================
% TITEL
% ============================================================================
\noindent
\begin{minipage}[t]{0.80\textwidth}
    \vspace{0pt}
    {\Large\bfseries Bohr-Vertiefung}\\[0.3em]
    {\normalsize Ionisierung \(\cdot\) \"Uberg\"ange \(\cdot\) Z-Erweiterung \(\cdot\) Fraunhofer \(\cdot\) Grenzen}\\[0.4em]
    Name: \luecke{5cm} \hfill Datum: \luecke{2.5cm}
\end{minipage}
\hfill
\begin{minipage}[t]{0.16\textwidth}
    \vspace{0pt}
    \raggedleft
    \href{https://devbox42.github.io/sciencesim/physik/kl12-GK/01-photoeffekt/07-DS-bohr-vertiefung/LP-07-bohr-vertiefung.html}{\qrcode[height=1.4cm]{https://devbox42.github.io/sciencesim/physik/kl12-GK/01-photoeffekt/07-DS-bohr-vertiefung/LP-07-bohr-vertiefung.html}}
\end{minipage}

\vspace{0.2em}
\begin{hinweisbox}
\centering
\textbf{Scanne den QR-Code und bearbeite den Lernpfad. F\"ulle dieses Blatt parallel aus.}
\end{hinweisbox}

\vspace{0.5em}

% ============================================================================
% AUFGABE 1: IONISIERUNGSENERGIE
% ============================================================================
\aufgabe{Aufgabe 1: Ionisierungsenergie} \sterne{2}

\vspace{0.3em}
\begin{merkbox}
\textbf{Ionisierungsenergie} = Energie, um ein Elektron \textbf{komplett} aus dem Atom zu entfernen.

\vspace{0.3em}
\begin{itemize}
    \item Grundzustand: $E_1 = $ \luecke{2cm}
    \item Komplett entfernt: $E_\infty = $ \luecke{2cm}
    \item Ionisierungsenergie: $E_\text{ion} = E_\infty - E_1 = $ \luecke{3cm}
\end{itemize}

\vspace{0.3em}
Die Ionisierungsenergie von Wasserstoff betr\"agt: $E_\text{ion} = $ \luecke{2cm}
\end{merkbox}

\vspace{0.3em}
\textbf{a)} Berechne die Ionisierungsenergie eines Elektrons, das sich auf $n = 3$ befindet.

\vspace{1.8cm}

$E_\text{ion}(n=3) = $ \luecke{3cm}

\vspace{0.5em}

% ============================================================================
% ENERGIENIVEAUSCHEMA (RP-Pflicht: skaliert)
% ============================================================================
\aufgabe{Energieniveauschema Wasserstoff} \sterne{2}

\vspace{0.3em}
Zeichne die fehlenden \"Uberg\"ange als Pfeile ein und beschrifte sie mit $\Delta E$ und $\lambda$.

\vspace{0.3em}
\begin{center}
\begin{tikzpicture}[scale=1, >=latex]
    % Achse
    \draw[->] (0,-0.3) -- (0,7.5) node[above]{\small $E$ (eV)};

    % Energieniveaus (skaliert: 0 eV = 7cm, -13.6 eV = 0cm)
    % Skalierung: y = (E + 13.6) / 13.6 * 7
    % E1 = -13.6 -> y = 0
    % E2 = -3.40 -> y = 5.25
    % E3 = -1.51 -> y = 6.22
    % E4 = -0.85 -> y = 6.56
    % E_inf = 0   -> y = 7.0

    % Niveaulinien
    \draw[thick] (0.5,0) -- (5,0) node[right]{\small $n=1$: $E_1 = \SI{-13,6}{eV}$};
    \draw[thick] (0.5,5.25) -- (5,5.25) node[right]{\small $n=2$: $E_2 = \SI{-3,40}{eV}$};
    \draw[thick] (0.5,6.22) -- (5,6.22) node[right]{\small $n=3$: $E_3 = $ \luecke{1.5cm}};
    \draw[thick] (0.5,6.56) -- (5,6.56) node[right]{\small $n=4$: $E_4 = $ \luecke{1.5cm}};
    \draw[thick, dashed] (0.5,7.0) -- (5,7.0) node[right]{\small $n=\infty$: $E_\infty = \SI{0}{eV}$ (Ionisierung)};

    % Beispiel-Uebergang: n=3 -> n=2 (vorgegeben)
    \draw[thick, physik, ->] (1.5,6.22) -- (1.5,5.25);
    \node[left, physik, font=\footnotesize] at (1.2,5.7) {H$_\alpha$};
    \node[left, font=\scriptsize] at (1.2,5.35) {\SI{656}{nm}};

    % Leere Pfeile zum Einzeichnen (gestrichelt als Hilfe)
    \draw[thin, gray, dashed, ->] (2.5,6.56) -- (2.5,5.25);
    \node[left, gray, font=\scriptsize] at (2.3,5.9) {?};

    \draw[thin, gray, dashed, ->] (3.5,6.22) -- (3.5,0);
    \node[left, gray, font=\scriptsize] at (3.3,3.1) {?};

    % Ionisierungspfeil
    \draw[thick, red, ->] (4.5,0) -- (4.5,7.0);
    \node[right, red, font=\scriptsize] at (4.6,3.5) {$E_\text{ion}$};

\end{tikzpicture}
\end{center}

\vspace{0.3em}
Zeichne ein: $n=4 \to 2$ (H$_\beta$) und $n=3 \to 1$ (Lyman-$\beta$). Beschrifte mit $\lambda$.

\vspace{0.5em}

% ============================================================================
% AUFGABE 2: \"UBERGANG BERECHNEN
% ============================================================================
\aufgabe{Aufgabe 2: \"Ubergangsrechnung Wasserstoff} \sterne{2}

\vspace{0.3em}
\begin{formelbox}
\textbf{Schnellformel:} \quad $\lambda = \dfrac{1240\;\text{eV\,nm}}{\Delta E}$ \qquad mit \quad $\Delta E = |E_{n_1} - E_{n_2}|$

\vspace{0.3em}
\textbf{Farbtabelle:} violett \SI{380}{}-\SI{430}{nm} | blau \SI{430}{}-\SI{470}{} | cyan \SI{470}{}-\SI{500}{} | gr\"un \SI{500}{}-\SI{560}{} | gelb \SI{560}{}-\SI{590}{} | orange \SI{590}{}-\SI{630}{} | rot \SI{630}{}-\SI{700}{nm}
\end{formelbox}

\vspace{0.3em}
\textbf{a)} Berechne Wellenl\"ange und Farbe f\"ur den \"Ubergang $n = 6 \to n = 2$.

\vspace{0.2em}
\begin{tabular}{@{}p{4.5cm}p{10cm}@{}}
$E_6 = -13{,}6 / 36 = $ & \luecke{3cm} \\[0.3cm]
$\Delta E = |E_6 - E_2| = $ & \luecke{4cm} \\[0.3cm]
$\lambda = 1240 / \Delta E = $ & \luecke{3cm} \\[0.3cm]
Farbe: & \luecke{3cm} \\
\end{tabular}

\vspace{0.5em}
\textbf{b)} Vervollst\"andige die Tabelle der Balmer-Serie (alle \"Uberg\"ange auf $n = 2$):

\vspace{0.3em}
\renewcommand{\arraystretch}{1.4}
\begin{tabular}{|c|c|c|c|c|}
\hline
\textbf{\"Ubergang} & $\mathbf{E_{n_1}}$ (eV) & $\boldsymbol{\Delta E}$ (eV) & $\boldsymbol{\lambda}$ (nm) & \textbf{Farbe} \\
\hline
$3 \to 2$ & $-1{,}51$ & $1{,}89$ & $656$ & rot \\
\hline
$4 \to 2$ & $-0{,}85$ & & & \\
\hline
$5 \to 2$ & $-0{,}544$ & & & \\
\hline
$6 \to 2$ & & & & \\
\hline
\end{tabular}

\vspace{0.5em}

% ============================================================================
% AUFGABE 3: ABSORPTION
% ============================================================================
\aufgabe{Aufgabe 3: Absorption} \sterne{3}

\vspace{0.3em}
Licht mit $\lambda = \SI{434}{nm}$ trifft auf Wasserstoff im Grundzustand.

\vspace{0.2em}
\textbf{a)} Berechne die Energie des Photons in eV.

\vspace{1.5cm}

$E_\text{Photon} = $ \luecke{3cm}

\vspace{0.3em}
\textbf{b)} Pr\"ufe: Kann dieses Photon von einem H-Atom im Grundzustand ($n=1$) absorbiert werden? Begr\"unde.

\textit{Tipp: Absorption ist nur m\"oglich, wenn $E_\text{Photon}$ genau der Differenz zweier Energieniveaus entspricht.}

\vspace{2cm}

\newpage

% ============================================================================
% AUFGABE 4: Z-ERWEITERUNG
% ============================================================================
\aufgabe{Aufgabe 4: Wasserstoff-\"ahnliche Ionen (Z-Erweiterung)} \sterne{3}

\vspace{0.3em}
\begin{merkbox}
F\"ur Ionen mit \textbf{einem Elektron} und Kernladungszahl $Z$:
$$E_n = -13{,}6\;\text{eV} \cdot \frac{Z^2}{n^2}$$

Beispiele: H ($Z=1$), He$^+$ ($Z=2$), Li$^{2+}$ ($Z=3$)
\end{merkbox}

\vspace{0.3em}
\textbf{a)} Berechne die Energieniveaus $E_1$, $E_2$ und $E_3$ f\"ur He$^+$ ($Z = 2$).

\vspace{0.2em}
\begin{tabular}{@{}p{3cm}p{11cm}@{}}
$E_1 = $ & \luecke{6cm} \\[0.4cm]
$E_2 = $ & \luecke{6cm} \\[0.4cm]
$E_3 = $ & \luecke{6cm} \\
\end{tabular}

\vspace{0.5em}
\textbf{b)} Berechne f\"ur He$^+$ den \"Ubergang $n = 4 \to n = 3$: Wellenl\"ange und Spektralbereich.

\vspace{0.2em}
\begin{tabular}{@{}p{4.5cm}p{10cm}@{}}
$E_4 = $ & \luecke{4cm} \\[0.3cm]
$\Delta E = $ & \luecke{4cm} \\[0.3cm]
$\lambda = $ & \luecke{4cm} \\[0.3cm]
Spektralbereich: & \luecke{4cm} (sichtbar / UV / IR) \\
\end{tabular}

\vspace{0.5em}
\textbf{c)} Erkl\"are, warum bei He$^+$ die Energiedifferenzen gr\"osser sind als bei H.

\vspace{1.8cm}

% ============================================================================
% AUFGABE 5: FRAUNHOFER-LINIEN
% ============================================================================
\aufgabe{Aufgabe 5: Fraunhofer-Linien und Sonnenspektrum} \sterne{2}

\vspace{0.3em}
\begin{merkbox}
\textbf{Absorptionsspektrum:} Die k\"uhlere Gasschicht der Sonnenatmosph\"are absorbiert bestimmte Wellenl\"angen aus dem kontinuierlichen Spektrum. Es entstehen \textbf{dunkle Linien} (Fraunhofer-Linien).

\vspace{0.3em}
\textbf{Prinzip:} Jedes Element absorbiert genau die Wellenl\"angen, die es auch \luecke{3cm} w\"urde.
\end{merkbox}

\vspace{0.3em}
\textbf{a)} Im Sonnenspektrum findet man dunkle Linien bei $\lambda = \SI{656}{nm}$ und $\lambda = \SI{486}{nm}$.

Welchem Element und welchen \"Uberg\"angen entsprechen diese Linien?

\vspace{0.3em}
\begin{tabular}{|c|c|c|c|}
\hline
$\boldsymbol{\lambda}$ & \textbf{Fraunhofer-Linie} & \textbf{Element} & \textbf{\"Ubergang} \\
\hline
\SI{656}{nm} & C & & \\
\hline
\SI{486}{nm} & F & & \\
\hline
\end{tabular}

\vspace{0.5em}
\textbf{b)} Im Spektrum eines Sterns fehlt die Natrium-D-Linie bei \SI{589}{nm}. In einem anderen Stern ist sie vorhanden.

Was kann man \"uber die chemische Zusammensetzung der beiden Sternatmosph\"aren schliessen?

\vspace{1.5cm}

\textbf{c)} \textbf{Natrium-Niveauschema:} Bei Na h\"angen die Energieniveaus auch vom Drehimpuls $l$ ab (nicht nur von $n$).

\vspace{0.3em}
\begin{center}
\begin{tikzpicture}[scale=0.9, >=latex]
    % Achse
    \draw[->] (0,-0.3) -- (0,5.8) node[above]{\small $E$ (eV)};

    % Niveaulinien (y = E + 5.14, skaliert 1:1 in cm)
    \draw[thick] (0.5,0) -- (4.5,0) node[right]{\small 3s: $E = \SI{-5,14}{eV}$ (Grundzustand)};
    \draw[thick] (0.5,2.10) -- (4.5,2.10) node[right]{\small 3p: $E = \SI{-3,04}{eV}$};
    \draw[thick] (0.5,3.19) -- (4.5,3.19) node[right]{\small 4s: $E = \SI{-1,95}{eV}$};
    \draw[thick] (0.5,3.62) -- (4.5,3.62) node[right]{\small 3d: $E = \SI{-1,52}{eV}$};
    \draw[thick, dashed] (0.5,5.14) -- (4.5,5.14) node[right]{\small Ionisierung: $\SI{0}{eV}$};

    % Gestrichelter Pfeil als Hinweis
    \draw[thin, gray, dashed, ->] (2.5,2.10) -- (2.5,0);
    \node[left, gray, font=\scriptsize] at (2.3,1.05) {D-Linie?};

\end{tikzpicture}
\end{center}

\vspace{0.2em}
Zeichne den \"Ubergang f\"ur die gelbe D-Linie ($\lambda = \SI{589}{nm}$) als Pfeil ein und pr\"ufe:

$\Delta E = 1240 / 589 = $ \luecke{2cm} eV \quad Passt das zu $|E_\text{3p} - E_\text{3s}|$? \luecke{2cm}

% ============================================================================
% AUFGABE 6: GRENZEN DES BOHR-MODELLS
% ============================================================================
\aufgabe{Aufgabe 6: St\"arken und Grenzen des Bohr-Modells} \sterne{2}

\vspace{0.3em}
Ordne die Aussagen zu:

\vspace{0.3em}
\renewcommand{\arraystretch}{1.5}
\begin{tabular}{|p{7cm}|c|c|}
\hline
\textbf{Aussage} & \textbf{St\"arke} & \textbf{Grenze} \\
\hline
Erkl\"art Spektrallinien von Wasserstoff & $\square$ & $\square$ \\
\hline
Versagt bei Helium (2 Elektronen) & $\square$ & $\square$ \\
\hline
Erkl\"art Ionisierungsenergie & $\square$ & $\square$ \\
\hline
Keine Aussage \"uber Intensit\"aten & $\square$ & $\square$ \\
\hline
Absorption und Emission erkl\"arbar & $\square$ & $\square$ \\
\hline
Feinstruktur nicht erkl\"arbar & $\square$ & $\square$ \\
\hline
\end{tabular}

\vspace{0.5em}
\textbf{Erg\"anze:} Die moderne Quantenmechanik ersetzt die festen Bahnen durch \luecke{3cm},

also dreidimensionale \luecke{5cm}.

\vspace{0.5em}

% ============================================================================
% SELBSTCHECK
% ============================================================================
\aufgabe{Selbstcheck: Klausurvorbereitung}

\vspace{0.3em}
\renewcommand{\arraystretch}{1.3}
\begin{tabular}{|p{10.5cm}|c|c|c|}
\hline
& \textbf{?} & \textbf{OK} & \textbf{!!} \\
\hline
Photogleichung anwenden ($E_\text{kin} = hf - W_A$) & $\square$ & $\square$ & $\square$ \\
\hline
Einsteinsche Gerade: Steigung = $h$, y-Achsenabschnitt = $-W_A$ & $\square$ & $\square$ & $\square$ \\
\hline
De-Broglie-Wellenl\"ange berechnen ($\lambda = h/p$) & $\square$ & $\square$ & $\square$ \\
\hline
Welle-Teilchen-Dualismus erkl\"aren & $\square$ & $\square$ & $\square$ \\
\hline
Drei Bohr-Postulate formulieren & $\square$ & $\square$ & $\square$ \\
\hline
\"Uberg\"ange berechnen: $\Delta E \to \lambda \to$ Farbe & $\square$ & $\square$ & $\square$ \\
\hline
Z-Erweiterung anwenden (He$^+$, Li$^{2+}$) & $\square$ & $\square$ & $\square$ \\
\hline
Fraunhofer-Linien erkl\"aren (Absorption) & $\square$ & $\square$ & $\square$ \\
\hline
Ionisierungsenergie bestimmen & $\square$ & $\square$ & $\square$ \\
\hline
Grenzen des Bohr-Modells benennen & $\square$ & $\square$ & $\square$ \\
\hline
\end{tabular}

\vspace{0.4em}
\begin{hinweisbox}
\textbf{Klausur: Fr 27.02.2026} | 60 BE | 90 min | Themen: Photoeffekt, de Broglie, Doppelspalt, Bohr, Spektralanalyse

\vspace{0.2em}
Bereite dich \"uber die Winterferien gezielt auf offene Punkte (?) vor!
\end{hinweisbox}

\end{document}
